%! TEX program = xelatex
%! TEX encoding = UTF-8 Unicode

\documentclass[internationalmaster]{hquThesis}
\usepackage{lineno}
\usepackage{amsmath,amsfonts,amssymb,textcomp,ulem}
\usepackage{booktabs,tabularray,float}
\graphicspath{{figures/}}

\hypersetup{colorlinks=true,linkcolor=black,citecolor=black}

\begin{document}
\makecover

% \include{decision}
% \include{declaration}

\frontmatter{
    \include{abstract}
    \tableofcontents
}

\mainmatter



\chapter{绪论}

\section{选题背景及意义}

板结构应用背景、屈曲/振动失效、Mindlin板剪切自锁问题、研究意义

\section{国内外研究历史及现状}

按方法分类综述（缩减积分、混合元、无网格法、有限元无网格混合离散），每段含方法溯源、核心思想、优缺点

\section{本文主要内容}

(1)(2)(3)列出主要贡献 + 各章内容安排

\chapter{剪切自锁问题}

本章阐明剪切自锁现象的成因与抑制途径，为后续混合离散格式的构造提供理论依据。首先从 Mindlin 板理论的位移场与应变-位移关系出发，分析薄板极限下剪切应变约束导致的数值刚化机制；其次归纳传统有限元中减缩积分与位移-剪力混合离散两类免自锁方案的原理与局限；随后引入约束比与 LBB 条件，说明最优约束比的选取原则；最后介绍本文采用的无网格增强免自锁方案的基本思想与 Mindlin 板运动学方程。

\section{剪切自锁现象的物理机制}

剪切自锁物理机制的详细推导与论述

\section{传统有限元免剪切自锁的方案}

\subsection{减缩积分方案}

减缩积分方案的原理与算例

\subsection{位移-剪力混合离散方案}

位移-剪力混合离散方案的详细论述

\section{最优约束比与无网格增强免自锁方案}

LBB 条件与约束比的数学推导

\section{Mindlin 板运动学方程}

考虑 Mindlin 板的三维位移场 $\bar{u}_i(x_1,x_2,x_3)$：
\begin{equation}
\begin{cases}
\bar{u}_\alpha(x_1,x_2,x_3) = u_\alpha(x_1,x_2) - x_3\varphi_\alpha(x_1,x_2), & \alpha = 1,2 \\
\bar{u}_3(x_1,x_2,x_3) = w(x_1,x_2),
\end{cases}
\label{eq:mindlin-displacement}
\end{equation}
其中 $u_\alpha$ 为中面面内位移，$w$ 为横向挠度，$\varphi_\alpha$ 为法线纤维的转角，$i,j=1,2,3$，$\alpha,\beta=1,2$，采用爱因斯坦求和约定。在小变形假设下，线性应变张量为 $\varepsilon_{ij} = \frac{1}{2}(u_{i,j}+u_{j,i})$。

面内应变分量（$\alpha,\beta=1,2$）为
\begin{equation}
\varepsilon_{\alpha\beta} = \frac{1}{2}(u_{\alpha,\beta}+u_{\beta,\alpha})
= \frac{x_3}{2}(\varphi_{\alpha,\beta}+\varphi_{\beta,\alpha})
\equiv x_3 \kappa_{\alpha\beta},
\label{eq:inplane-strain}
\end{equation}
其中 $\kappa_{\alpha\beta} = \frac{1}{2}(\varphi_{\alpha,\beta}+\varphi_{\beta,\alpha})$ 为中面曲率。横向剪切应变（$\alpha=1,2$）为
\begin{equation}
\gamma_\alpha \equiv 2\varepsilon_{\alpha 3} = \bar{u}_{\alpha,3}+\bar{u}_{3,\alpha}
= w_{,\alpha} - \varphi_\alpha.
\label{eq:transverse-shear-strain}
\end{equation}

假设板处于平面应力状态（$\sigma_{33}=0$），材料为各向同性线弹性。面内应力与应变关系为 $\sigma_{\alpha\beta}=C_{\alpha\beta\gamma\delta}\varepsilon_{\gamma\delta}$，四阶弹性张量为：
\begin{equation}
C_{\alpha\beta\gamma\delta} = \frac{E}{1-\nu^2}\left[\nu\delta_{\alpha\beta}\delta_{\gamma\delta} + \frac{1-\nu}{2}(\delta_{\alpha\gamma}\delta_{\beta\delta}+\delta_{\alpha\delta}\delta_{\beta\gamma})\right],
\label{eq:elastic-tensor}
\end{equation}
其中 $E$ 为杨氏模量，$\nu$ 为泊松比。横向剪应力需引入剪力修正系数 $\kappa$，有
\begin{equation}
\sigma_{\alpha 3} = \kappa G \gamma_\alpha, \qquad G = \frac{E}{2(1+\nu)},
\label{eq:shear-constitutive}
\end{equation}
$\kappa$ 通常取 $5/6$。

\section{小结}



\chapter{自由振动剪切自锁有限元无网格混合离散分析}
% 第三章：自由振动分析公式推导
% 包含：混合离散动力学方程、能量方程、弱形式
% [修订v1] 补 \label/\eqref 交叉引用；统一 \varphi（原 \psi）；质量子块改 m_{IJ}（原 M_{IJ}，与弯矩合力 M_{\alpha\beta} 撞符号）

本章针对 Mindlin 中厚板的自由振动问题，基于第二章建立的运动学方程，推导混合离散格式下的动力学控制方程与广义特征值问题。首先推导板的动能变分与应变能变分，通过哈密顿原理建立自由振动的虚功方程；其次进行离散化，得到材料刚度矩阵与质量矩阵的各子块表达式；最后讨论混合离散格式的稳定性条件。

\section{混合离散动力学方程与特征值问题}

考虑 Mindlin 中厚板的自由振动问题，求解目标为结构的固有频率 $\omega$ 和固有振型。自由振动对应的广义特征值问题为：
\begin{equation}
(\mathbf{K} - \omega^2 \mathbf{M})\mathbf{D} = \mathbf{0}
\label{eq:vib-eigen}
\end{equation}
式中，$\mathbf{K}$ 为弹性刚度矩阵，$\mathbf{M}$ 为质量矩阵，$\mathbf{D}$ 为节点自由度向量，$\omega$ 为固有圆频率。式\eqref{eq:vib-eigen}即为后续离散化工作的求解目标。

\section{振动分析能量方程}

由式\eqref{eq:mindlin-displacement}给出的位移场、式\eqref{eq:inplane-strain}与式\eqref{eq:transverse-shear-strain}的应变-位移关系以及式\eqref{eq:shear-constitutive}的材料本构可知，本章在第二章的基础上进行动力分析，首先推导板的动能变分与应变能变分，再通过哈密顿原理建立自由振动的虚功方程。板的动能变分为：
\begin{equation}
\delta T = \int_V \rho \ddot{u}_i \delta u_i\, dV
\label{eq:vib-kinetic-variation}
\end{equation}
将式\eqref{eq:mindlin-displacement}代入，利用中面位置 $x_3=0$ 的对称性以及密度均匀假设，定义惯性系数
\begin{equation}
I_0 = \int_{-h/2}^{h/2} \rho\, dx_3 = \rho h, \qquad I_2 = \int_{-h/2}^{h/2} \rho x_3^2\, dx_3 = \frac{\rho h^3}{12}
\label{eq:inertia-coefficients}
\end{equation}
式中，$I_0$ 为平动惯性系数，$I_2$ 为转动惯性系数，$h$ 为板厚。中面对称使 $\int_{-h/2}^{h/2} \rho x_3 dx_3 = 0$，可得动能变分的中面积分形式：
\begin{equation}
\delta T = \int_\Omega \left(I_0 \ddot{w}\,\delta w + I_2 \ddot{\varphi}_\alpha \delta\varphi_\alpha\right) d\Omega
\label{eq:vib-kinetic-midplane}
\end{equation}
式中，$\Omega$ 为板的中面，$\ddot{w}$ 为横向加速度，$\ddot{\varphi}_\alpha$ 为法线转角加速度。

板的应变能变分为：
\begin{equation}
\delta U = \int_V \sigma_{ij} \delta\varepsilon_{ij}\, dV
\label{eq:vib-strain-variation}
\end{equation}
内虚功密度 $\sigma_{ij}\delta\varepsilon_{ij}$ 包含面内分量与横向分量（$\sigma_{33}=0$ 不计）：
\begin{equation}
\sigma_{ij}\delta\varepsilon_{ij} = \sigma_{\alpha\beta}\delta\varepsilon_{\alpha\beta} + 2\sigma_{\alpha 3}\delta\varepsilon_{\alpha 3}
\label{eq:virtual-work-density}
\end{equation}
利用 Mindlin 位移场对应的应变变分 $\delta\varepsilon_{\alpha\beta} = x_3\delta\kappa_{\alpha\beta}$，$\delta\varepsilon_{\alpha 3} = \frac{1}{2}\delta\gamma_\alpha$，代入并对厚度积分，定义弯矩合力与剪力合力
\begin{equation}
M_{\alpha\beta} \equiv \int_{-h/2}^{h/2} x_3 \sigma_{\alpha\beta}\, dx_3 = \frac{h^3}{12} C_{\alpha\beta\gamma\delta}\kappa_{\gamma\delta}, \qquad Q_\alpha \equiv \int_{-h/2}^{h/2} \sigma_{\alpha 3}\, dx_3 = \kappa G h \gamma_\alpha
\label{eq:stress-resultants}
\end{equation}
式中，$M_{\alpha\beta}$ 为弯矩合力，$Q_\alpha$ 为剪力合力，$C_{\alpha\beta\gamma\delta}$ 为式\eqref{eq:elastic-tensor}定义的四阶弹性张量，$\kappa$ 为剪切修正系数。将式\eqref{eq:stress-resultants}代入式\eqref{eq:vib-strain-variation}并对厚度积分，内虚功可写为
\begin{equation}
\delta U = \int_\Omega \left(M_{\alpha\beta}\delta\kappa_{\alpha\beta} + Q_\alpha \delta\gamma_\alpha\right) d\Omega
\label{eq:vib-internal-virtual-work}
\end{equation}

根据哈密顿原理，系统在任意时间区间 $[t_1, t_2]$ 内的动能变分与应变能变分满足：
\begin{equation}
\int_{t_1}^{t_2} (\delta T - \delta U)\, dt = 0,
\label{eq:hamilton-principle}
\end{equation}
对任意虚位移 $\delta u_i$，可得弹性动力学的虚功方程：
\begin{equation}
\int_V (\sigma_{ij}\delta\varepsilon_{ij} + \rho \ddot{u}_i \delta u_i)\, dV = 0.
\label{eq:vib-virtual-work}
\end{equation}
板的体积积分可分解为中面积分与厚度积分：$\int_V dV = \int_\Omega \int_{-h/2}^{h/2} dx_3\, d\Omega$，其中 $\Omega$ 为板的中面。将式\eqref{eq:vib-internal-virtual-work}与式\eqref{eq:vib-kinetic-midplane}代入式\eqref{eq:vib-virtual-work}，得到 Mindlin 板自由振动的虚功方程：
\begin{equation}
\int_\Omega \left[M_{\alpha\beta}\delta\kappa_{\alpha\beta} + Q_\alpha \delta\gamma_\alpha - I_0 \ddot{w}\,\delta w - I_2 \ddot{\varphi}_\alpha \delta \varphi_\alpha\right] d\Omega = 0.
\label{eq:vib-plate-virtual-work}
\end{equation}
式中，前两项为弹性内力虚功（弯曲与剪切），后两项为惯性力虚功。式\eqref{eq:vib-plate-virtual-work}即为自由振动问题离散化的出发点。

\section{弱形式}

对式\eqref{eq:vib-plate-virtual-work}各变量取变分并令其为零，将离散场代入后得到广义特征值方程。位移场和转角场的形函数插值为：
\begin{equation}
w(\mathbf{x}) = N_I w_I, \qquad \varphi_\alpha(\mathbf{x}) = N_I \varphi_{\alpha I}
\label{eq:shape-function-interpolation}
\end{equation}
式中，$N_I$ 为节点 $I$ 的形函数，$w_I$ 为节点横向挠度，$\varphi_{\alpha I}$ 为节点法线转角。对空间坐标求偏导，微分算子直接作用在形函数上：$w_{,\alpha} = N_{I,\alpha}w_I$，$\varphi_{\alpha,\beta} = N_{I,\beta}\varphi_{\alpha I}$。虚位移做同阶离散近似。

将式\eqref{eq:shape-function-interpolation}代入式\eqref{eq:vib-internal-virtual-work}，材料刚度矩阵各子块为
\begin{equation}
K_{IJ}^{ww} = \int_\Omega \kappa G h\, N_{I,\alpha} N_{J,\alpha}\, d\Omega, \qquad K_{I,J\gamma}^{w\varphi} = -\int_\Omega \kappa G h\, N_{I,\gamma} N_J\, d\Omega
\label{eq:stiffness-ww-wphi}
\end{equation}
\begin{equation}
K_{I\alpha,J\gamma}^{\varphi\varphi} = \int_\Omega \left(N_{I,\beta} D_{\alpha\beta\gamma\eta} N_{J,\eta} + \kappa G h\, N_I N_J \delta_{\alpha\gamma}\right) d\Omega
\label{eq:stiffness-phiphi}
\end{equation}
式中，$K_{IJ}^{ww}$ 为横向挠度对应的剪切刚度子块，$K_{I,J\gamma}^{w\varphi}$ 为挠度与转角耦合的刚度子块，$K_{I\alpha,J\gamma}^{\varphi\varphi}$ 为转角对应的弯曲与剪切刚度子块，$D_{\alpha\beta\gamma\eta} = \frac{h^3}{12}C_{\alpha\beta\gamma\eta}$ 为弯曲刚度张量，$\delta_{\alpha\gamma}$ 为 Kronecker 符号。

质量矩阵展开后的对称矩阵形式为
\begin{equation}
\mathbf{M} = \begin{bmatrix}
\mathbf{M}^{ww} & \mathbf{0} & \mathbf{0} \\
\mathbf{0} & \mathbf{M}^{\varphi_1\varphi_1} & \mathbf{0} \\
\mathbf{0} & \mathbf{0} & \mathbf{M}^{\varphi_2\varphi_2}
\end{bmatrix}
\label{eq:mass-matrix-block}
\end{equation}
其中横向挠度对应的质量子块为：
\begin{equation}
m_{IJ}^{ww} = \int_\Omega I_0 N_I N_J\, d\Omega = \int_\Omega \rho h\, N_I N_J\, d\Omega
\label{eq:mass-ww}
\end{equation}
转角自由度对应的质量子块为：
\begin{equation}
m_{I\alpha,J\beta}^{\varphi\varphi} = \int_\Omega I_2 \delta_{\alpha\beta} N_I N_J\, d\Omega = \int_\Omega \frac{\rho h^3}{12} \delta_{\alpha\beta} N_I N_J\, d\Omega
\label{eq:mass-phiphi}
\end{equation}
式中，$m_{IJ}^{ww}$ 为横向平动质量子块，$m_{I\alpha,J\beta}^{\varphi\varphi}$ 为转动质量子块，$\delta_{\alpha\beta}$ 为 Kronecker 符号。中面面内位移与横向挠度不耦合，质量矩阵呈块对角形式。

线性问题中面内与弯曲不耦合，全局特征值方程呈块对角结构：
\begin{equation}
\begin{bmatrix}
\mathbf{K}_{ww} & \mathbf{K}_{w\varphi} \\
\mathbf{K}_{\varphi w} & \mathbf{K}_{\varphi\varphi}
\end{bmatrix}
\begin{bmatrix}
\mathbf{W} \\
\boldsymbol{\Psi}
\end{bmatrix}
= \omega^2
\begin{bmatrix}
\mathbf{M}_{ww} & \mathbf{0} \\
\mathbf{0} & \mathbf{M}_{\varphi\varphi}
\end{bmatrix}
\begin{bmatrix}
\mathbf{W} \\
\boldsymbol{\Psi}
\end{bmatrix}
\label{eq:vib-global-eigen}
\end{equation}
式中，$\mathbf{W}$ 为节点挠度向量，$\boldsymbol{\Psi}$ 为节点转角向量，$\mathbf{K}_{ww}$、$\mathbf{K}_{w\varphi}$、$\mathbf{K}_{\varphi\varphi}$ 分别为由式\eqref{eq:stiffness-ww-wphi}与式\eqref{eq:stiffness-phiphi}组集而成的全局刚度子矩阵。此即 Mindlin 板自由振动的全局离散特征值方程。求解广义特征值问题式\eqref{eq:vib-eigen}，可得到自然频率 $\omega$ 和对应模态。

振动问题中质量矩阵 $\mathbf{M}$ 为对称正定，自由振动特征值问题的适定性要求混合离散格式满足 LBB 稳定性条件。本文采用的剪切应力场插值阶数经约束比分析验证，可保证 inf-sup 常数有正下界，从而在薄板极限下避免剪切自锁并维持刚度矩阵的正定性。当板厚趋近于零时，最低阶频率收敛于 Kirchhoff 薄板理论的解析解。

% 方板自由振动数值算例
% 格式对齐第四章：三线表（无竖线）、caption在表上方、不用\scriptsize、图片\linewidth

\section{数值算例}

通过数值算例系统检验不同边界条件下，混合格式在自由振动分析中的计算精度。本文数值算例采用无量纲材料参数，弹性模量 $E=1.0$，泊松比 $\nu=0.3$，密度 $\rho=1.0$，板厚 $h=10^{-3}$，板边长 $a=b=1.0$，剪切修正系数 $\kappa=5/6$。

无量纲频率参数定义为
\begin{equation}
\lambda = \frac{\bar{\omega}}{\pi^2}, \qquad \bar{\omega} = \omega a^2 \sqrt{\frac{\rho h}{D^b}}, \qquad D^b = \frac{E h^3}{12(1-\nu^2)}.
\end{equation}

参考解取自 Liew 等人\cite{liew1993} 采用 $pb$-$2$ Rayleigh-Ritz 法计算的结果，对应 $a/b=1.0$、$t/b=0.001$。

\subsection{频率对比}

五种数值方法用于对比：标准有限元（FEM）、缩减积分有限元（FEM缩减积分）、三变量混合有限元（FEM三变量）、无网格法（MF）、改进无网格法（MF2）。网格密度从 $20\times20$ 到 $30\times30$。

% TODO[数据口径]：当前FEM列为h=1e-2结果，待重跑为h=1e-3后更新
% TODO[方法说明]：MF与MF2最大相对差仅2.7e-4，非独立方法，需在表注说明
% TODO[对称性]：CSCS与SCSC在方板自由振动下由90°旋转对称严格等价，需合并或注明

\begin{table}[H]
\centering
\caption{CCCC边界前十阶无量纲频率 $\lambda$ 对比（Tri3网格，$25\times25$，$t/b=0.001$）}
\label{tab:vib_cccc}
\begin{tabular}{crrrrrrr}
\hline
阶次 & 标准解\cite{liew1993} & FEM & FEM缩减积分 & FEM三变量 & MF & MF2 \\
\hline
1 & 3.6463 & 9.0161 & 81.6867 & 3.6755 & 3.6466 & 3.6466 \\
2 & 7.4364 & 16.7062 & 148.2861 & 7.5235 & 7.4383 & 7.4385 \\
3 & 7.4364 & 19.7881 & 181.7857 & 7.5248 & 7.4383 & 7.4386 \\
4 & 10.9647 & 25.6354 & 226.2437 & 11.1483 & 10.9676 & 10.9682 \\
5 & 13.3315 & 32.0603 & 288.7722 & 13.5552 & 13.3378 & 13.3378 \\
6 & 13.3951 & 34.7997 & 321.2210 & 13.6267 & 13.4009 & 13.4009 \\
7 & 16.7177 & 36.6469 & 321.2819 & 17.0684 & 16.7241 & 16.7245 \\
8 & 16.7177 & 44.0710 & 396.6453 & 17.0803 & 16.7242 & 16.7248 \\
9 & 21.3296 & 49.6168 & 433.3729 & 21.8123 & 21.3440 & 21.3437 \\
10 & 21.3296 & 52.0191 & 471.5575 & 21.8188 & 21.3441 & 21.3438 \\
\hline
\end{tabular}
\end{table}

\begin{table}[H]
\centering
\caption{SSSS边界前十阶无量纲频率 $\lambda$ 对比（Tri3网格，$25\times25$，$t/b=0.001$）}
\label{tab:vib_ssss}
\begin{tabular}{crrrrrrr}
\hline
阶次 & 标准解 & FEM & FEM缩减积分 & FEM三变量 & MF & MF2 \\
\hline
1 & 2.0000 & 5.2493 & 47.5683 & 2.0078 & 2.0000 & 2.0000 \\
2 & 5.0000 & 11.2975 & 99.6490 & 5.0177 & 5.0006 & 5.0008 \\
3 & 5.0000 & 13.9179 & 128.0485 & 5.0292 & 5.0006 & 5.0008 \\
4 & 8.0000 & 18.9447 & 165.9420 & 8.0589 & 8.0002 & 8.0009 \\
5 & 10.0000 & 24.3030 & 218.3133 & 10.0673 & 10.0040 & 10.0042 \\
6 & 10.0000 & 26.7288 & 247.1585 & 10.0778 & 10.0040 & 10.0042 \\
7 & 13.0000 & 28.5535 & 248.4308 & 13.1272 & 13.0019 & 13.0029 \\
8 & 13.0000 & 35.2044 & 315.4769 & 13.1377 & 13.0021 & 13.0031 \\
9 & 17.0000 & 40.1562 & 348.1101 & 17.1846 & 17.0135 & 17.0139 \\
10 & 17.0000 & 41.8796 & 379.1305 & 17.1972 & 17.0136 & 17.0141 \\
\hline
\end{tabular}
\end{table}

\begin{table}[H]
\centering
\caption{CSCS边界前十阶无量纲频率 $\lambda$ 对比（Tri3网格，$25\times25$，$t/b=0.001$）}
\label{tab:vib_cscs}
\begin{tabular}{crrrrrrr}
\hline
阶次 & 标准解\cite{liew1993} & FEM & FEM缩减积分 & FEM三变量 & MF & MF2 \\
\hline
1 & 2.9334 & 7.3545 & 66.6355 & 2.9543 & 2.9336 & 2.9336 \\
2 & 5.5465 & 13.4000 & 119.3601 & 5.5886 & 5.5476 & 5.5478 \\
3 & 7.0242 & 17.7497 & 162.2562 & 7.1026 & 7.0257 & 7.0260 \\
4 & 9.5833 & 22.0418 & 193.7378 & 9.7214 & 9.5848 & 9.5855 \\
5 & 10.3564 & 26.9177 & 244.5636 & 10.4559 & 10.3613 & 10.3614 \\
6 & 13.0797 & 32.2880 & 282.2141 & 13.2967 & 13.0851 & 13.0852 \\
7 & 14.2053 & 32.6540 & 299.0921 & 14.4385 & 14.2097 & 14.2105 \\
8 & 15.6815 & 39.7010 & 355.5860 & 15.9922 & 15.6854 & 15.6862 \\
9 & 17.2597 & 44.5333 & 388.0839 & 17.4838 & 17.2743 & 17.2746 \\
10 & 20.2442 & 44.5796 & 408.3700 & 20.7280 & 20.2487 & 20.2499 \\
\hline
\end{tabular}
\end{table}

% TODO[对称性]：SCSC与CSCS在方板自由振动下由90°旋转对称严格等价，数据完全相同，可考虑合并

\begin{table}[H]
\centering
\caption{FSCS边界前十阶无量纲频率 $\lambda$ 对比（Tri3网格，$25\times25$，$t/b=0.001$）}
\label{tab:vib_fscs}
\begin{tabular}{crrrrrrr}
\hline
阶次 & 标准解\cite{liew1993} & FEM & FEM缩减积分 & FEM三变量 & MF & MF2 \\
\hline
1 & 1.2853 & 3.4320 & 31.1699 & 1.2888 & 1.2857 & 1.2860 \\
2 & 3.3502 & 8.2615 & 72.5739 & 3.3689 & 3.3492 & 3.3495 \\
3 & 4.2253 & 11.1267 & 102.6841 & 4.2428 & 4.2276 & 4.2293 \\
4 & 6.3846 & 15.0915 & 131.8066 & 6.4329 & 6.3852 & 6.3867 \\
5 & 7.3352 & 19.4851 & 176.4565 & 7.3935 & 7.3344 & 7.3349 \\
6 & 9.1817 & 22.6276 & 201.0210 & 9.2476 & 9.1893 & 9.1925 \\
7 & 10.4521 & 24.2808 & 218.1133 & 10.5718 & 10.4502 & 10.4519 \\
8 & 11.3372 & 29.9415 & 268.3376 & 11.4564 & 11.3437 & 11.3466 \\
9 & 13.3161 & 34.0026 & 296.0380 & 13.4776 & 13.3175 & 13.3180 \\
10 & 15.4787 & 35.0810 & 319.2567 & 15.7094 & 15.4811 & 15.4843 \\
\hline
\end{tabular}
\end{table}

\begin{table}[H]
\centering
\caption{FSSS边界前十阶无量纲频率 $\lambda$ 对比（Tri3网格，$25\times25$，$t/b=0.001$）}
\label{tab:vib_fsss}
\begin{tabular}{crrrrrrr}
\hline
阶次 & 标准解\cite{liew1993} & FEM & FEM缩减积分 & FEM三变量 & MF & MF2 \\
\hline
1 & 1.1847 & 3.1356 & 28.4668 & 1.1865 & 1.1841 & 1.1844 \\
2 & 2.8127 & 7.2761 & 64.2912 & 2.8232 & 2.8116 & 2.8119 \\
3 & 4.1743 & 10.6714 & 97.7653 & 4.1899 & 4.1763 & 4.1778 \\
4 & 5.9846 & 13.7398 & 119.9629 & 6.0196 & 5.9853 & 5.9868 \\
5 & 6.2679 & 17.6729 & 161.0025 & 6.2981 & 6.2669 & 6.2673 \\
6 & 9.1486 & 21.6411 & 188.6070 & 9.2128 & 9.1569 & 9.1600 \\
7 & 9.5730 & 23.2016 & 210.6866 & 9.6495 & 9.5716 & 9.5733 \\
8 & 11.0356 & 28.0091 & 250.4310 & 11.1336 & 11.0418 & 11.0447 \\
9 & 11.7212 & 31.6099 & 277.9959 & 11.8096 & 11.7231 & 11.7237 \\
10 & 14.7558 & 32.8124 & 296.9094 & 14.9263 & 14.7586 & 14.7617 \\
\hline
\end{tabular}
\end{table}

\subsection{收敛分析}

收敛分析采用绝对差值
\begin{equation}
e = |\lambda_{\mathrm{num}} - \lambda_{\mathrm{ref}}|,
\end{equation}
其中 $\lambda_{\mathrm{num}}$ 为数值解，$\lambda_{\mathrm{ref}}$ 为标准解。图\ref{vib_convergence} 给出六种边界条件下第一阶模态的收敛曲线。图中横坐标为 $\log_{10}(h)$，$h=1/n$ 为单元尺寸；纵坐标为 $\log_{10}(e)$，误差值越往下表示结果越精确。

\begin{figure}[H]
\centering
\includegraphics[width=0.9\textwidth]{vibration/convergence_first_mode_all.png}
\caption{六种边界条件第一阶模态无量纲频率的对数误差收敛曲线}
\label{vib_convergence}
\end{figure}

% TODO[图件中文化+去边框]：当前收敛图标签为英文且有坐标轴边框，需重画为无框图

\subsection{模态对比}

四边简支（SSSS）方板的自由振动模态具有解析形式，可直接作为振型对比的基准。图\ref{vib_modal_ssss} 给出 SSSS 边界下前六阶模态的变形图对比。每张图的列为不同数值方法，行为模态阶次。

% TODO[图件中文化+去边框]：当前模态图标签为英文且有坐标轴边框，需重画

\begin{figure}[H]
\centering
\textbf{(a) 第 1 阶}\\[1mm]
\includegraphics[width=\linewidth]{vibration/modal_shapes_SSSS.png}\\[3mm]
\caption{四边简支（SSSS）方板前六阶模态变形图对比}
\label{vib_modal_ssss}
\end{figure}

% TODO[待补充]：CCCC等其他边界的模态图（当前无标准模态，待后续补充）
% TODO[待补充]：收敛率定量值（MF、FEM三变量混合元对各边界的网格收敛斜率）
% TODO[待补充]：结果分析文字（读图三件套：图例解读→读法说明→结论）

\section{小结}

% TODO: 第三章小结

\chapter{屈曲分析剪切自锁有限元无网格混合离散分析}
% 第四章：屈曲分析公式推导
% 包含：混合离散屈曲控制方程、能量方程、弱形式
% [修订v1] 1) 补 \label/\eqref 交叉引用；2) §4.3 子块推导合并为「引导句+一组公式+式中」，删除 5 处 \boxed 与逐式重复；
%          3) 补变分符号 δ（原式 ∫σε = δW_ext 不成立）；4) 修正 D_{\alpha\beta\gamma\eta}、\bar\sigma_{\alpha\beta}、N_{J,\gamma} 等下标笔误；
%          5) 统一 \varphi（原 \phi 混用）与 \mathbf{K}_b/\mathbf{K}_g 记号。

本章针对 Mindlin 中厚板的屈曲问题，基于第二章建立的运动学方程，推导混合离散格式下的屈曲控制方程与广义特征值问题。首先采用 Green-Lagrange 应变张量推导屈曲问题的虚功方程；其次进行离散化，得到材料刚度矩阵与几何刚度矩阵的各子块表达式；最后讨论混合离散格式的稳定性条件。

\section{混合离散屈曲控制方程与特征值问题}

考虑 Mindlin 中厚板的屈曲问题。板在面内压力作用下，当荷载增大至临界值时，原始平面内平衡状态失去稳定性，板突然发生横向挠度（屈曲分岔）。与强度失效不同，屈曲失效通常发生在应力远低于材料屈服强度时，具有突发性。求解目标为临界荷载系数 $\lambda$ 和屈曲模态，对应的广义特征值问题为：
\begin{equation}
(\mathbf{K} + \lambda \mathbf{K}_G)\boldsymbol{\phi} = \mathbf{0}
\label{eq:buck-eigen}
\end{equation}
式中，$\mathbf{K}$ 为弹性刚度矩阵，$\mathbf{K}_G$ 为几何刚度矩阵，$\lambda$ 为屈曲荷载因子，$\boldsymbol{\phi}$ 为屈曲模态。临界屈曲荷载为 $P_{cr} = \lambda_{cr} P_{ref}$，其中 $P_{ref}$ 为参考预加荷载，$\lambda_{cr}$ 为最小正特征值。

面内位移变分 $u_\alpha$ 不会与预应力度 $\bar{\sigma}_{\alpha\beta}$ 产生导致屈曲的几何力矩，因此几何刚度矩阵中与 $u_\alpha$ 对应的行列皆为零。实际控制板分岔屈曲的广义特征值问题子块矩阵方程为：
\begin{equation}
(\mathbf{K}_b + \lambda \mathbf{K}_g)\mathbf{b} = \mathbf{0}
\label{eq:buck-reduced-eigen}
\end{equation}
式中，$\mathbf{K}_b$ 为面外材料刚度矩阵，$\mathbf{K}_g$ 为几何刚度矩阵，$\mathbf{b} = [\,w,\ \varphi_1,\ \varphi_2\,]^{\mathrm{T}}$ 为面外节点自由度向量。式\eqref{eq:buck-reduced-eigen}即本章离散化工作的求解目标。

\section{屈曲分析能量方程}

由第二章建立的 Mindlin 板位移场可知，屈曲分析需考虑几何非线性，采用完整的 Green-Lagrange 应变张量：
\begin{equation}
\varepsilon_{ij} = \frac{1}{2}(\bar{u}_{i,j} + \bar{u}_{j,i} + \bar{u}_{k,i}\bar{u}_{k,j})
\label{eq:green-lagrange-strain}
\end{equation}

将式\eqref{eq:mindlin-displacement}代入，面内应变分量 $\varepsilon_{\alpha\beta}$ 按 $x_3$ 的幂次分项归类为：
\begin{equation}
\begin{aligned}
\varepsilon_{\alpha\beta} ={}& \underbrace{\frac{1}{2}(u_{\alpha,\beta} + u_{\beta,\alpha} + u_{\gamma,\alpha}u_{\gamma,\beta})}_{\text{薄膜应变}} + \underbrace{\frac{1}{2}w_{,\alpha}w_{,\beta}}_{\text{von K\'arm\'an 项}} \\
&- x_3\underbrace{\frac{1}{2}(\varphi_{\alpha,\beta} + \varphi_{\beta,\alpha} + u_{\gamma,\alpha}\varphi_{\gamma,\beta} + \varphi_{\gamma,\alpha}u_{\gamma,\beta})}_{\text{广义曲率}} \\
&+ x_3^2\underbrace{\frac{1}{2}\varphi_{\gamma,\alpha}\varphi_{\gamma,\beta}}_{\text{高阶旋转项}}
\end{aligned}
\label{eq:strain-power-expansion}
\end{equation}

横向剪切应变 $\varepsilon_{\alpha 3}$ 和厚度方向法向应变 $\varepsilon_{33}$ 分别为：
\begin{equation}
\varepsilon_{\alpha 3} = \frac{1}{2}(w_{,\alpha} - \varphi_\alpha - u_{\beta,\alpha}\varphi_\beta) + \frac{x_3}{2}\varphi_{\beta,\alpha}\varphi_\beta, \qquad \varepsilon_{33} = \frac{1}{2}\varphi_\alpha\varphi_\alpha
\label{eq:shear-normal-strain}
\end{equation}
式中，$u_\alpha$、$w$、$\varphi_\alpha$ 为式\eqref{eq:mindlin-displacement}定义的中面位移与转角，$x_3$ 为厚度方向坐标。上述各应变分量的完整逐项展开过程参见附录\ref{app:strain-expansion}。

根据虚功原理，板的内虚功等于外虚功：
\begin{equation}
\int_\Omega\int_{-h/2}^{h/2} \sigma_{ij}\delta\varepsilon_{ij}\, dx_3\, d\Omega = \delta W_{ext}
\label{eq:virtual-work-principle}
\end{equation}
将应变张量按面内分量、横向剪切分量和厚度方向分量拆分为三部分积分：
\begin{equation}
\int_\Omega\int_{-h/2}^{h/2} \sigma_{\alpha\beta}\delta\varepsilon_{\alpha\beta}\, dx_3 d\Omega + 2\int_\Omega\int_{-h/2}^{h/2} \sigma_{\alpha 3}\delta\varepsilon_{\alpha 3}\, dx_3 d\Omega + \int_\Omega\int_{-h/2}^{h/2} \sigma_{33}\delta\varepsilon_{33}\, dx_3 d\Omega = \delta W_{ext}
\label{eq:virtual-work-split}
\end{equation}

第一项为面内应变项。将总应力 $\sigma_{\alpha\beta}$ 拆分为预应力度与屈曲增量应力：
\begin{equation}
\sigma_{\alpha\beta} = \bar{\sigma}_{\alpha\beta} + \sigma_{\alpha\beta}^{inc}
\label{eq:stress-decomposition}
\end{equation}
其中，预应力度 $\bar{\sigma}_{\alpha\beta}$ 沿厚度均匀分布（与 $x_3$ 无关）；增量应力 $\sigma_{\alpha\beta}^{inc}$ 随厚度弹性变化。将式\eqref{eq:strain-power-expansion}代入并完全展开，面内应变分量与应力的乘积共含 8 项：(1) 式\eqref{eq:strain-power-expansion}中薄膜应变部分与预应力度 $\bar{\sigma}_{\alpha\beta}$ 的乘积；(2) 薄膜应变部分与增量应力 $\sigma_{\alpha\beta}^{inc}$ 的乘积；(3) 广义曲率项 $-x_3\kappa_{\alpha\beta}$ 与预应力度的乘积；(4) 广义曲率项与增量应力的乘积；(5) von K\'arm\'an 项 $\frac{1}{2}w_{,\alpha}w_{,\beta}$ 与预应力度的乘积；(6) von K\'arm\'an 项与增量应力的乘积；(7) 高阶旋转项 $x_3^2\eta_{\alpha\beta}$ 与预应力度的乘积；(8) 高阶旋转项与增量应力的乘积。依据小变形假设和屈曲分岔点的特性，第 1、2 项涉及膜应变，其一阶变分为零；第 6、8 项为增量应力与二阶位移梯度的乘积，属三阶小量舍去；第 3 项因预应力度沿厚度均匀分布与 $x_3$ 无关，$\int_{-h/2}^{h/2} x_3 dx_3 = 0$ 恒消去。保留第 4、5、7 项，对厚度积分后定义弯矩合力：
\begin{equation}
M_{\alpha\beta} = \int_{-h/2}^{h/2} x_3 \sigma_{\alpha\beta}^{inc}\, dx_3
\label{eq:moment-resultant-inc}
\end{equation}
式中，$M_{\alpha\beta}$ 为屈曲增量弯矩合力。第一项积分的最终结果为：
\begin{equation}
\int_\Omega\int_{-h/2}^{h/2} \sigma_{\alpha\beta}\delta\varepsilon_{\alpha\beta}\, dx_3 d\Omega = \int_\Omega \delta\kappa_{\alpha\beta}M_{\alpha\beta}\, d\Omega + \int_\Omega h\bar{\sigma}_{\alpha\beta}w_{,\alpha}\delta w_{,\beta}\, d\Omega + \int_\Omega \frac{h^3}{12}\bar{\sigma}_{\alpha\beta}\varphi_{\gamma,\alpha}\delta\varphi_{\gamma,\beta}\, d\Omega
\label{eq:inplane-term-result}
\end{equation}
式中，右端第一项为屈曲增量弯矩对应的弯曲虚功，第二项与第三项为预应力度对应的几何虚功（分别为挠度梯度项与转角梯度项）。

第二项为横向剪切项。对于纯面内加载的板，预屈曲横向剪应力 $\bar{\sigma}_{\alpha 3} = 0$，故 $\sigma_{\alpha 3} = \sigma_{\alpha 3}^{inc}$。由虎克定律，引入剪切修正系数 $\kappa$ 后：
\begin{equation}
\sigma_{\alpha 3}^{inc} = \kappa G(w_{,\alpha} - \varphi_\alpha)
\label{eq:shear-stress-inc}
\end{equation}
代入并对厚度积分（$\int_{-h/2}^{h/2} dx_3 = h$），注意工程剪切应变 $\gamma_\alpha = 2\varepsilon_{\alpha 3}$，第二项的积分结果为：
\begin{equation}
2\int_\Omega\int_{-h/2}^{h/2} \sigma_{\alpha 3}\delta\varepsilon_{\alpha 3}\, dx_3 d\Omega = \int_\Omega \kappa G h\,(w_{,\alpha} - \varphi_\alpha)\delta(w_{,\alpha} - \varphi_\alpha)\, d\Omega
\label{eq:shear-term-result}
\end{equation}

第三项为厚度方向法向应变项，在平面应力假设 $\sigma_{33}=0$ 下贡献为零。

将三项合并，得到屈曲问题的完整虚功方程：
\begin{equation}
\begin{aligned}
\delta W_{ext} ={}& \int_\Omega \delta\kappa_{\alpha\beta}M_{\alpha\beta}\, d\Omega + \int_\Omega \kappa G h\,(w_{,\alpha} - \varphi_\alpha)\delta(w_{,\alpha} - \varphi_\alpha)\, d\Omega \\
&+ \int_\Omega h\bar{\sigma}_{\alpha\beta} w_{,\alpha}\delta w_{,\beta}\, d\Omega + \int_\Omega \frac{h^3}{12}\bar{\sigma}_{\alpha\beta} \varphi_{\gamma,\alpha}\delta\varphi_{\gamma,\beta}\, d\Omega
\end{aligned}
\label{eq:buck-virtual-work-complete}
\end{equation}
式中，前两项为弹性内力虚功（弯曲与剪切），后两项为预应力度对应的几何刚度虚功。当几何刚度项的绝对值等于弹性刚度项时，系统总刚度矩阵奇异，达到临界屈曲状态。式\eqref{eq:buck-virtual-work-complete}即为屈曲问题离散化的出发点。

\section{弱形式}

对式\eqref{eq:buck-virtual-work-complete}各变量取变分并令其为零，将离散场代入后得到离散形式的广义特征值方程。位移场和转角场的形函数插值为：
\begin{equation}
w(\mathbf{x}) = N_I w_I, \qquad \varphi_\alpha(\mathbf{x}) = N_I \varphi_{\alpha I}
\label{eq:buck-shape-function-interpolation}
\end{equation}
式中，$N_I$ 为节点 $I$ 的形函数，$w_I$ 为节点横向挠度，$\varphi_{\alpha I}$ 为节点法线转角。对空间坐标求偏导，微分算子直接作用在形函数上：$w_{,\alpha} = N_{I,\alpha}w_I$，$\varphi_{\alpha,\beta} = N_{I,\beta}\varphi_{\alpha I}$。虚位移做同阶离散近似。

材料刚度矩阵子块展开如下：面外材料刚度矩阵 $\mathbf{K}_b$ 由弯曲刚度（Bending）与横向剪切刚度（Shear）共同组成。对应自由度 $[\,w,\ \varphi_1,\ \varphi_2\,]^{\mathrm{T}}$，展开后的对称矩阵形式为
\begin{equation}
\mathbf{K}_b = \begin{bmatrix}
K_{ww} & K_{w\varphi_1} & K_{w\varphi_2} \\
K_{\varphi_1 w} & K_{\varphi_1\varphi_1} & K_{\varphi_1\varphi_2} \\
K_{\varphi_2 w} & K_{\varphi_2\varphi_1} & K_{\varphi_2\varphi_2}
\end{bmatrix}
\label{eq:kb-matrix}
\end{equation}
式中，对角块 $K_{ww}$、$K_{\varphi_1\varphi_1}$、$K_{\varphi_2\varphi_2}$ 为同一自由度内部的刚度贡献，非对角块为不同自由度之间的耦合刚度贡献，矩阵关于主对角线对称。

以 $K_{ww}$ 为例说明离散化过程。将式\eqref{eq:buck-shape-function-interpolation}代入式\eqref{eq:buck-virtual-work-complete}的横向剪切项，注意 $w_{,\alpha} = N_{I,\alpha}w_I$、$\delta w_{,\alpha} = N_{J,\alpha}\delta w_J$，有
\begin{equation}
\begin{aligned}
\int_\Omega \kappa G h\, w_{,\alpha}\delta w_{,\alpha}\, d\Omega
&= \int_\Omega \kappa G h\, (N_{I,\alpha}w_I)(N_{J,\alpha}\delta w_J)\, d\Omega \\
&= \delta w_J \left[\int_\Omega \kappa G h\, N_{I,\alpha}N_{J,\alpha}\, d\Omega\right] w_I
\end{aligned}
\label{eq:kww-chain}
\end{equation}
式中，方括号内的积分即为挠度自由度对应的剪切刚度子块 $K_{IJ}^{ww} = \int_\Omega \kappa G h\, N_{I,\alpha}N_{J,\alpha}\, d\Omega$。其余子块作相同处理，将弯曲项与剪切项分别代入并整理，各子块的积分表达式汇总为
\begin{equation}
\begin{aligned}
K_{IJ}^{ww} &= \int_\Omega \kappa G h\, N_{I,\alpha}N_{J,\alpha}\, d\Omega, &
K_{I,J\gamma}^{w\varphi} &= -\int_\Omega \kappa G h\, N_{I,\gamma}N_J\, d\Omega, \\
K_{I\gamma,J\gamma}^{\varphi\varphi,s} &= \int_\Omega \kappa G h\, N_I N_J\, d\Omega, &
K_{I\alpha,J\gamma}^{\varphi\varphi,b} &= \int_\Omega N_{I,\beta}D_{\alpha\beta\gamma\eta}N_{J,\eta}\, d\Omega
\end{aligned}
\label{eq:kb-subblocks}
\end{equation}
式中，$K_{IJ}^{ww}$ 为横向剪切对挠度的贡献，$K_{I,J\gamma}^{w\varphi}$ 为挠度与转角的耦合贡献，$K_{I\gamma,J\gamma}^{\varphi\varphi,s}$ 为横向剪切对转角的贡献，$K_{I\alpha,J\gamma}^{\varphi\varphi,b}$ 为弯曲对转角的贡献，上标 $s$ 与 $b$ 分别表示剪切（Shear）与弯曲（Bending）；$D_{\alpha\beta\gamma\eta} = \frac{h^3}{12}C_{\alpha\beta\gamma\eta}$ 为弯曲刚度张量，$C_{\alpha\beta\gamma\eta}$ 为式\eqref{eq:elastic-tensor}定义的四阶弹性张量；$\kappa$ 为剪切修正系数，薄板或中厚板通常取 $5/6$。

几何刚度矩阵子块展开如下：由式\eqref{eq:buck-virtual-work-complete}后两项可知，几何刚度仅与挠度梯度和转角梯度有关，故其子块矩阵同样对应自由度 $[\,w,\ \varphi_1,\ \varphi_2\,]^{\mathrm{T}}$，展开后的对称矩阵形式为
\begin{equation}
\mathbf{K}_g = \begin{bmatrix}
K_g^{ww} & & \\
& K_g^{\varphi_1\varphi_1} & K_g^{\varphi_1\varphi_2} \\
& K_g^{\varphi_2\varphi_1} & K_g^{\varphi_2\varphi_2}
\end{bmatrix}
\label{eq:kg-matrix}
\end{equation}
式中，空白位置对应的子块为零，即预应力度不产生挠度与转角之间的耦合几何刚度。

以 $K_g^{ww}$ 为例，将式\eqref{eq:buck-shape-function-interpolation}代入式\eqref{eq:buck-virtual-work-complete}的几何虚功项，有
\begin{equation}
\begin{aligned}
\int_\Omega h\bar{\sigma}_{\alpha\beta} w_{,\alpha}\delta w_{,\beta}\, d\Omega
&= \int_\Omega h\bar{\sigma}_{\alpha\beta}(N_{I,\alpha}w_I)(N_{J,\beta}\delta w_J)\, d\Omega \\
&= \delta w_J \left[\int_\Omega h\bar{\sigma}_{\alpha\beta}N_{I,\alpha}N_{J,\beta}\, d\Omega\right] w_I
\end{aligned}
\label{eq:kgww-chain}
\end{equation}
式中，方括号内的积分为挠度自由度对应的几何刚度子块 $K_{g,IJ}^{ww} = \int_\Omega h\bar{\sigma}_{\alpha\beta}N_{I,\alpha}N_{J,\beta}\, d\Omega$。转角自由度作相同处理，各子块的积分表达式汇总为
\begin{equation}
\begin{aligned}
K_{g,IJ}^{ww} &= \int_\Omega h\bar{\sigma}_{\alpha\beta}N_{I,\alpha}N_{J,\beta}\, d\Omega, &
K_{g,I\gamma,J\gamma}^{\varphi\varphi} &= \int_\Omega \frac{h^3}{12}\bar{\sigma}_{\alpha\beta}N_{I,\alpha}N_{J,\beta}\, d\Omega = \frac{h^2}{12}K_{g,IJ}^{ww}
\end{aligned}
\label{eq:kg-subblocks}
\end{equation}
式中，$K_{g,IJ}^{ww}$ 为挠度自由度对应的几何刚度子块，$K_{g,I\gamma,J\gamma}^{\varphi\varphi}$ 为转角自由度对应的几何刚度子块；因 $\int_{-h/2}^{h/2}x_3^2 dx_3 = h^3/12$，转角子块与挠度子块之比恒为 $h^2/12$，二者无需分别积分。

将上述所有子块合并，针对中厚板面外屈曲模式，最终的广义特征值矩阵方程为：
\begin{equation}
\left(
\begin{bmatrix}
K_{ww} & K_{w\varphi_1} & K_{w\varphi_2} \\
K_{\varphi_1 w} & K_{\varphi_1\varphi_1} & K_{\varphi_1\varphi_2} \\
K_{\varphi_2 w} & K_{\varphi_2\varphi_1} & K_{\varphi_2\varphi_2}
\end{bmatrix}
+ \lambda
\begin{bmatrix}
K_g^{ww} & & \\
& \frac{h^2}{12}K_g^{ww} & \\
& & \frac{h^2}{12}K_g^{ww}
\end{bmatrix}
\right)
\begin{bmatrix} w \\ \varphi_1 \\ \varphi_2 \end{bmatrix}
= \begin{bmatrix} 0 \\ 0 \\ 0 \end{bmatrix}
\label{eq:buck-global-eigen}
\end{equation}
式中，$\lambda$ 为待求的屈曲荷载因子，$w$ 与 $\varphi_1$、$\varphi_2$ 为节点自由度向量，几何刚度矩阵的非零块由式\eqref{eq:kg-subblocks}给出。式\eqref{eq:buck-global-eigen}与式\eqref{eq:buck-reduced-eigen}形式一致，求解该广义特征值问题即可得到临界荷载系数与屈曲模态。

% TODO: 待补充——混合离散格式的LBB稳定性分析
% 如需补充，应结合第二章的约束比理论，说明本格式的剪切应力场插值阶数满足inf-sup条件
% 可参考徐洋涛论文第3章的写法：独立成节，含特征值问题推导和约束比数值验证

% 方板屈曲数值算例
% 包含：CCCC, SSSS, CSCS, FSCS, FSSS, SCSC 六种边界条件
% [修订v1] 1) 修正事实性错误：v01–v12 为第 1–12 阶屈曲模态编号，不是厚跨比；
%          2) 补齐图表叙述（参数表、对比表、收敛图的读图与结论）；
%          3) 繁体「模態」→「模态」；补 \label 与 \cite；
%          4) 所有算例数据与相对误差均取自 date\ 目录收敛序列末值（ndiv=35），核验脚本见
%             D:\Joker\deepseek_harness\scripts\data_check\report_convergence.py，报告见 D:\Joker\refs\屈曲收敛数据核验.md
% [待办-数据] 表 CH4_compare_table 的「FEM缩减积分」列与「FEM」列在各边界下末值完全相同（序列逐值一致），
%             疑为同一份数据重复导出，暂不足以支撑「缩减积分」作为独立方法的对比结论，需重算或删列。

\section{数值算例}

针对单轴均匀压缩载荷，系统评估混合离散格式在极薄极限下的临界载荷计算能力。本节首先给出方板计算模型、材料与几何参数以及验证基准；其次在六种边界条件下对比各方法的临界荷载系数及其网格收敛性；最后给出第 1 至第 12 阶屈曲模态下各方法的对比图形。

\subsection{计算模型与参数}

方板求解域为 $\Omega = [-0.5, 0.5] \times [-0.5, 0.5]$，采用均匀节点离散。四条边分别编号为 1（底边）、2（右边）、3（顶边）、4（左边），通过组合不同的边界约束条件（固支 C、简支 S、自由 F）构成六种典型工况：CCCC、SSSS、CSCS、SCSC、FSCS 与 FSSS。图\ref{CH4_plate_model} 给出方板计算模型与边界条件示意图，图中边界箭头表示各边的位移约束方向。

\begin{figure}[H]
    \centering
    \includegraphics[width=0.7\textwidth]{plate_model_generated.png}
    \caption{方板计算模型与边界条件示意图}
    \label{CH4_plate_model}
\end{figure}

本文数值算例采用无量纲材料参数，具体取值列于表\ref{CH4_params} 中。方板边长 $a=b=1.0$，板厚 $h=10^{-2}$，厚跨比 $h/a=1/100$，属于薄板范畴。剪切修正系数取 Mindlin 板理论常用值 $\kappa=5/6$。

\begin{table}[H]
\centering
\caption{方板屈曲算例材料与几何参数}
\label{CH4_params}
\begin{tabular}{lll}
\hline
\textbf{参数} & \textbf{符号} & \textbf{取值} \\
\hline
弹性模量 & $E$ & $1.0$ \\
泊松比 & $\nu$ & $0.3$ \\
密度 & $\rho$ & $1.0$ \\
板厚 & $h$ & $10^{-2}$ \\
板边长 & $a=b$ & $1.0$ \\
剪切修正系数 & $\kappa$ & $5/6$ \\
弯曲刚度 & $D^b$ & $\approx 9.16\times10^{-8}$ \\
剪切刚度 & $D^s$ & $\approx 3.21\times10^{-3}$ \\
预加应力 & $\sigma_{11},\sigma_{22},\sigma_{12}$ & $1.0,0,0$ \\
\hline
\end{tabular}
\end{table}

表\ref{CH4_params} 中，弹性模量、泊松比与密度取无量纲单位；预加应力 $\sigma_{11}=1.0$ 对应沿 $x$ 方向的单轴均匀压缩，$\sigma_{22}=\sigma_{12}=0$；弯曲刚度 $D^b = Eh^3/[12(1-\nu^2)]$ 与剪切刚度 $D^s = \kappa Gh$ 由几何与材料参数直接确定。表中剪切刚度比弯曲刚度高约 4 个数量级，薄板极限下横向剪切项成为约束条件的主要来源，剪切自锁即出现在该参数区间。

数值算例采用结构化网格离散，四边形单元（Quad4）网格规模为 $25\times25$ 至 $35\times35$，三角形单元（Tri3）网格规模对应约 $392$ 至 $800$ 个单元。典型网格如图\ref{CH4_mesh} 所示。

\begin{figure}[H]
    \centering
    \includegraphics[width=0.9\textwidth]{meshes/mesh_quad4_styled.png}
    \vspace{2mm}
    \includegraphics[width=0.9\textwidth]{meshes/mesh_tri3_styled.png}
    \caption{方板问题节点离散示意图}
    \label{CH4_mesh}
\end{figure}

\subsection{验证基准与收敛性分析}

屈曲临界荷载系数的验证基准基于 Kirchhoff 薄板理论。由于本算例厚跨比 $h/a=1/100$，属于薄板范畴，Mindlin 理论的计算结果应收敛于 Kirchhoff 理论的解析解或经典参考解。

对于四边简支（SSSS）方板，采用 Navier 双三角级数法可获得精确闭合解。设板沿 $x$ 方向承受单轴均布压力 $N_x$，屈曲控制方程为：
\begin{equation}
D\nabla^4 w + N_x \frac{\partial^2 w}{\partial x^2} = 0,
\label{eq:buck-kirchhoff-control}
\end{equation}
其中弯曲刚度 $D = Eh^3/[12(1-\nu^2)]$。取满足四边简支边界条件（$w=0,\ M_n=0$）的试函数：
\begin{equation}
w(x,y) = \sum_{m=1}^{\infty}\sum_{n=1}^{\infty} A_{mn} \sin\frac{m\pi x}{a} \sin\frac{n\pi y}{b},
\label{eq:navier-mode-shape}
\end{equation}
式中，$A_{mn}$ 为待定幅值系数，$m$、$n$ 分别为 $x$ 与 $y$ 方向的半波数。将式\eqref{eq:navier-mode-shape}代入式\eqref{eq:buck-kirchhoff-control}，可得方板最低阶（$m=n=1$）屈曲临界荷载系数：
\begin{equation}
k_{\mathrm{SSSS}} = \frac{(m^2+n^2)^2}{m^2} = 4.000.
\label{eq:ssss-buckling-coefficient}
\end{equation}

对于四边固支（CCCC）、混合边界（CSCS、SCSC）及含自由边（FSCS、FSSS）的方板，由于边界条件无法用简单三角级数自动满足，不存在精确闭合解，需通过 Levy 单三角级数法或 Rayleigh-Ritz 能量法求解超越方程获得数值参考解。本文采用 Timoshenko 与 Gere\cite{timoshenko1961} 及 Bui 等人\cite{bui2011} 给出的经典参考值作为验证基准，各边界条件的最低阶屈曲系数汇总于表\ref{CH4_compare_table}。

\begin{table}[H]
\centering
\caption{单轴受压方板屈曲临界荷载系数对比（三角形 Tri3 网格，$35\times35$，$t/b=0.01$）}
\label{CH4_compare_table}
\begin{tabular}{cccccccc}
\hline
边界条件 & 标准解\cite{timoshenko1961} & 文献\cite{bui2011} & MF & MF2 & FEM & FEM缩减积分 & FEM三变量 \\
\hline
CCCC & 10.0700 & 10.1080 & 10.0568 & 10.1040 & 8.2847 & 8.2847 & 10.0842 \\
SSSS & 4.0000 & 4.0050 & 3.9921 & 3.9987 & 3.5966 & 3.5966 & 4.0081 \\
CSCS & 7.6900 & 7.7020 & 7.6975 & 7.6617 & 6.2903 & 6.2903 & 7.7069 \\
SCSC & 6.7500 & 6.7800 & 6.7292 & 6.7620 & 5.8397 & 5.8397 & 6.7392 \\
FSCS & 1.7000 & 1.6930 & 1.6497 & 1.6528 & 1.5470 & 1.5470 & 1.6547 \\
FSSS & 1.4400 & 1.4390 & 1.3993 & 1.4019 & 1.3224 & 1.3224 & 1.4032 \\
\hline
\end{tabular}
\end{table}

从表中可以看出，六种边界条件下 MF 与 FEM 三变量混合元的屈曲临界荷载系数均接近标准解，而常规 FEM 的结果系统性偏低，且偏低幅度随边界约束增强而增大。四边固支工况下，标准解为 10.0700，MF 为 10.0568，FEM 三变量混合元为 10.0842，三者差值均在 0.02 以内；常规 FEM 仅为 8.2847，比标准解低 1.7853。四边简支工况的规律相同，MF 与 FEM 三变量混合元分别为 3.9921 与 4.0081，常规 FEM 为 3.5966，比标准解 4.0000 低 0.4034。含自由边的 FSCS 与 FSSS 工况下，常规 FEM 与 MF 的差值由四边固支的 1.7721 降至 0.1027 与 0.0769。自由边工况下板的屈曲模态以局部变形为主，弯曲与剪切的耦合程度低于全约束工况，剪切自锁的影响随之减弱。

在表\ref{CH4_compare_table} 的收敛序列基础上，图\ref{CH4_convergence} 给出 CCCC 与 SSSS 两种典型工况下各方法的误差随单元尺寸的变化。图中横坐标为 $\log_{10}(h)$，$h=1/(n-1)$ 为单元尺寸；纵坐标为 $\log_{10}(|K-K_{\mathrm{exact}}|/K_{\mathrm{exact}})$，误差值越往下表示结果越精确。图中第二个 MF 变体标注为 MF1，与表\ref{CH4_compare_table} 中的 MF2 为同一格式。

\begin{figure}[H]
    \centering
    \begin{tabular}{cc}
    \includegraphics[width=0.48\textwidth]{convergence/conv_tri3_CCCC.png} &
    \includegraphics[width=0.48\textwidth]{convergence/conv_tri3_SSSS.png} \\
    (a) 四边固支 CCCC & (b) 四边简支 SSSS
    \end{tabular}
    \caption{三角形（Tri3）网格下各方法屈曲临界荷载系数的对数误差收敛曲线}
    \label{CH4_convergence}
\end{figure}

从图中可以看出，MF 的误差曲线随单元尺寸减小单调下降，在双对数坐标下的斜率约为 $1.55$，四边简支工况的对应斜率约为 $1.26$，误差随网格加密稳定收敛。常规 FEM 的误差曲线呈相反趋势：网格由 $25\times25$ 加密至 $35\times35$ 时，其四边固支结果由 8.5316 下降至 8.2847，与标准解的差值由 1.5384 扩大至 1.7853，加密网格反而使结果进一步偏离薄板理论的解析解。常规 FEM 的误差不来自网格离散，而来自剪切约束项在薄板极限下的数值刚化，加密网格无法消除。FEM 三变量混合元的末值与 MF 相当，但其误差曲线在相邻网格间呈奇偶交替波动，波动幅度约为 0.04，尚不呈现单调收敛特征。

% [待办-数据] 上述收敛率由 date\ 目录序列拟合得到（斜率按 log10(err) 对 log10(h) 线性回归）。
% FSCS/FSSS 工况下 MF 的误差随网格加密基本不变（拟合斜率约 0.07），提示该两工况的参考值定义
% 或荷载归一化方式与其余工况不一致，需在定稿前核对基准来源。

\subsection{屈曲模态对比}

为考察各方法对高阶模态的解析能力，图\ref{CH4_buckling_CCCC_a} 至图\ref{CH4_buckling_SCSC_b} 给出六种边界条件下第 1 至第 12 阶屈曲模态的临界荷载系数对比结果。每幅图包含 6 组面板，每组对应一个屈曲模态阶次，组内自左至右依次为标准解、MF、MF1、FEM、FEM 缩减积分与 FEM 三变量混合元；方法标签仅标注于第 1 阶与第 7 阶模态所在的面板组中，图中 $v_{01}$ 至 $v_{12}$ 依次对应第 1 至第 12 阶屈曲模态。

四边固支方板的边界约束最强，其屈曲临界荷载系数在六种工况中最高，结果如图\ref{CH4_buckling_CCCC_a} 与图\ref{CH4_buckling_CCCC_b} 所示。

\begin{figure}[H]
    \centering
    \textbf{(a) 第 1 阶}\\[1mm]
    \includegraphics[width=\linewidth]{combined/layout_1x6_v9/CCCC_v01_1x6.png}\\[3mm]
    \textbf{(b) 第 2 阶}\\[1mm]
    \includegraphics[width=\linewidth]{combined/layout_1x6_v9/CCCC_v02_1x6.png}\\[3mm]
    \textbf{(c) 第 3 阶}\\[1mm]
    \includegraphics[width=\linewidth]{combined/layout_1x6_v9/CCCC_v03_1x6.png}\\[3mm]
    \textbf{(d) 第 4 阶}\\[1mm]
    \includegraphics[width=\linewidth]{combined/layout_1x6_v9/CCCC_v04_1x6.png}\\[3mm]
    \textbf{(e) 第 5 阶}\\[1mm]
    \includegraphics[width=\linewidth]{combined/layout_1x6_v9/CCCC_v05_1x6.png}\\[3mm]
    \textbf{(f) 第 6 阶}\\[1mm]
    \includegraphics[width=\linewidth]{combined/layout_1x6_v9/CCCC_v06_1x6.png}\\[3mm]
    \caption{四边固支（CCCC）方板第 1 至第 6 阶屈曲模态临界荷载系数对比}
    \label{CH4_buckling_CCCC_a}
\end{figure}
\begin{figure}[H]
    \centering
    \textbf{(a) 第 7 阶}\\[1mm]
    \includegraphics[width=\linewidth]{combined/layout_1x6_v9/CCCC_v07_1x6.png}\\[3mm]
    \textbf{(b) 第 8 阶}\\[1mm]
    \includegraphics[width=\linewidth]{combined/layout_1x6_v9/CCCC_v08_1x6.png}\\[3mm]
    \textbf{(c) 第 9 阶}\\[1mm]
    \includegraphics[width=\linewidth]{combined/layout_1x6_v9/CCCC_v09_1x6.png}\\[3mm]
    \textbf{(d) 第 10 阶}\\[1mm]
    \includegraphics[width=\linewidth]{combined/layout_1x6_v9/CCCC_v10_1x6.png}\\[3mm]
    \textbf{(e) 第 11 阶}\\[1mm]
    \includegraphics[width=\linewidth]{combined/layout_1x6_v9/CCCC_v11_1x6.png}\\[3mm]
    \textbf{(f) 第 12 阶}\\[1mm]
    \includegraphics[width=\linewidth]{combined/layout_1x6_v9/CCCC_v12_1x6.png}\\[3mm]
    \caption{四边固支（CCCC）方板第 7 至第 12 阶屈曲模态临界荷载系数对比}
    \label{CH4_buckling_CCCC_b}
\end{figure}

% TODO[图像核对] 补充本段观察句：第 1 阶面板中 MF、MF1 与 FEM 三变量混合元的振型是否与标准解一致，
% 常规 FEM 面板的振型相对标准解的偏差表现（波数、节线位置、幅值分布），以及该偏差随模态阶次升高的变化趋势。

四边简支方板的模态为式\eqref{eq:navier-mode-shape}的双三角级数形式，可直接作为振型对比的基准，结果如图\ref{CH4_buckling_SSSS_a} 与图\ref{CH4_buckling_SSSS_b} 所示。

\begin{figure}[H]
    \centering
    \textbf{(a) 第 1 阶}\\[1mm]
    \includegraphics[width=\linewidth]{combined/layout_1x6_v9/SSSS_v01_1x6.png}\\[3mm]
    \textbf{(b) 第 2 阶}\\[1mm]
    \includegraphics[width=\linewidth]{combined/layout_1x6_v9/SSSS_v02_1x6.png}\\[3mm]
    \textbf{(c) 第 3 阶}\\[1mm]
    \includegraphics[width=\linewidth]{combined/layout_1x6_v9/SSSS_v03_1x6.png}\\[3mm]
    \textbf{(d) 第 4 阶}\\[1mm]
    \includegraphics[width=\linewidth]{combined/layout_1x6_v9/SSSS_v04_1x6.png}\\[3mm]
    \textbf{(e) 第 5 阶}\\[1mm]
    \includegraphics[width=\linewidth]{combined/layout_1x6_v9/SSSS_v05_1x6.png}\\[3mm]
    \textbf{(f) 第 6 阶}\\[1mm]
    \includegraphics[width=\linewidth]{combined/layout_1x6_v9/SSSS_v06_1x6.png}\\[3mm]
    \caption{四边简支（SSSS）方板第 1 至第 6 阶屈曲模态临界荷载系数对比}
    \label{CH4_buckling_SSSS_a}
\end{figure}
\begin{figure}[H]
    \centering
    \textbf{(a) 第 7 阶}\\[1mm]
    \includegraphics[width=\linewidth]{combined/layout_1x6_v9/SSSS_v07_1x6.png}\\[3mm]
    \textbf{(b) 第 8 阶}\\[1mm]
    \includegraphics[width=\linewidth]{combined/layout_1x6_v9/SSSS_v08_1x6.png}\\[3mm]
    \textbf{(c) 第 9 阶}\\[1mm]
    \includegraphics[width=\linewidth]{combined/layout_1x6_v9/SSSS_v09_1x6.png}\\[3mm]
    \textbf{(d) 第 10 阶}\\[1mm]
    \includegraphics[width=\linewidth]{combined/layout_1x6_v9/SSSS_v10_1x6.png}\\[3mm]
    \textbf{(e) 第 11 阶}\\[1mm]
    \includegraphics[width=\linewidth]{combined/layout_1x6_v9/SSSS_v11_1x6.png}\\[3mm]
    \textbf{(f) 第 12 阶}\\[1mm]
    \includegraphics[width=\linewidth]{combined/layout_1x6_v9/SSSS_v12_1x6.png}\\[3mm]
    \caption{四边简支（SSSS）方板第 7 至第 12 阶屈曲模态临界荷载系数对比}
    \label{CH4_buckling_SSSS_b}
\end{figure}

% TODO[图像核对] 补充本段观察句：四边简支条件下模态重根现象（同一阶次出现成对相等或接近的临界荷载系数）
% 在各方法面板中的表现，以及 MF 与 FEM 三变量混合元对重根模态的分辨能力。

固支-简支交替（CSCS）工况的两组对边约束不同，其屈曲临界荷载系数介于四边固支与四边简支之间，结果如图\ref{CH4_buckling_CSCS_a} 与图\ref{CH4_buckling_CSCS_b} 所示。

\begin{figure}[H]
    \centering
    \textbf{(a) 第 1 阶}\\[1mm]
    \includegraphics[width=\linewidth]{combined/layout_1x6_v9/CSCS_v01_1x6.png}\\[3mm]
    \textbf{(b) 第 2 阶}\\[1mm]
    \includegraphics[width=\linewidth]{combined/layout_1x6_v9/CSCS_v02_1x6.png}\\[3mm]
    \textbf{(c) 第 3 阶}\\[1mm]
    \includegraphics[width=\linewidth]{combined/layout_1x6_v9/CSCS_v03_1x6.png}\\[3mm]
    \textbf{(d) 第 4 阶}\\[1mm]
    \includegraphics[width=\linewidth]{combined/layout_1x6_v9/CSCS_v04_1x6.png}\\[3mm]
    \textbf{(e) 第 5 阶}\\[1mm]
    \includegraphics[width=\linewidth]{combined/layout_1x6_v9/CSCS_v05_1x6.png}\\[3mm]
    \textbf{(f) 第 6 阶}\\[1mm]
    \includegraphics[width=\linewidth]{combined/layout_1x6_v9/CSCS_v06_1x6.png}\\[3mm]
    \caption{固支-简支交替（CSCS）方板第 1 至第 6 阶屈曲模态临界荷载系数对比}
    \label{CH4_buckling_CSCS_a}
\end{figure}
\begin{figure}[H]
    \centering
    \textbf{(a) 第 7 阶}\\[1mm]
    \includegraphics[width=\linewidth]{combined/layout_1x6_v9/CSCS_v07_1x6.png}\\[3mm]
    \textbf{(b) 第 8 阶}\\[1mm]
    \includegraphics[width=\linewidth]{combined/layout_1x6_v9/CSCS_v08_1x6.png}\\[3mm]
    \textbf{(c) 第 9 阶}\\[1mm]
    \includegraphics[width=\linewidth]{combined/layout_1x6_v9/CSCS_v09_1x6.png}\\[3mm]
    \textbf{(d) 第 10 阶}\\[1mm]
    \includegraphics[width=\linewidth]{combined/layout_1x6_v9/CSCS_v10_1x6.png}\\[3mm]
    \textbf{(e) 第 11 阶}\\[1mm]
    \includegraphics[width=\linewidth]{combined/layout_1x6_v9/CSCS_v11_1x6.png}\\[3mm]
    \textbf{(f) 第 12 阶}\\[1mm]
    \includegraphics[width=\linewidth]{combined/layout_1x6_v9/CSCS_v12_1x6.png}\\[3mm]
    \caption{固支-简支交替（CSCS）方板第 7 至第 12 阶屈曲模态临界荷载系数对比}
    \label{CH4_buckling_CSCS_b}
\end{figure}

% TODO[图像核对] 补充本段观察句：CSCS 工况下常规 FEM 面板相对标准解的偏差在低阶与高阶模态间是否变化，
% 以及 MF 面板对混合边界引起的模态局部化现象的再现程度。

自由-简支-固支-简支（FSCS）工况含一个自由边，边界约束弱于四边约束工况，结果如图\ref{CH4_buckling_FSCS_a} 与图\ref{CH4_buckling_FSCS_b} 所示。

\begin{figure}[H]
    \centering
    \textbf{(a) 第 1 阶}\\[1mm]
    \includegraphics[width=\linewidth]{combined/layout_1x6_v9/FSCS_v01_1x6.png}\\[3mm]
    \textbf{(b) 第 2 阶}\\[1mm]
    \includegraphics[width=\linewidth]{combined/layout_1x6_v9/FSCS_v02_1x6.png}\\[3mm]
    \textbf{(c) 第 3 阶}\\[1mm]
    \includegraphics[width=\linewidth]{combined/layout_1x6_v9/FSCS_v03_1x6.png}\\[3mm]
    \textbf{(d) 第 4 阶}\\[1mm]
    \includegraphics[width=\linewidth]{combined/layout_1x6_v9/FSCS_v04_1x6.png}\\[3mm]
    \textbf{(e) 第 5 阶}\\[1mm]
    \includegraphics[width=\linewidth]{combined/layout_1x6_v9/FSCS_v05_1x6.png}\\[3mm]
    \textbf{(f) 第 6 阶}\\[1mm]
    \includegraphics[width=\linewidth]{combined/layout_1x6_v9/FSCS_v06_1x6.png}\\[3mm]
    \caption{自由-简支-固支-简支（FSCS）方板第 1 至第 6 阶屈曲模态临界荷载系数对比}
    \label{CH4_buckling_FSCS_a}
\end{figure}
\begin{figure}[H]
    \centering
    \textbf{(a) 第 7 阶}\\[1mm]
    \includegraphics[width=\linewidth]{combined/layout_1x6_v9/FSCS_v07_1x6.png}\\[3mm]
    \textbf{(b) 第 8 阶}\\[1mm]
    \includegraphics[width=\linewidth]{combined/layout_1x6_v9/FSCS_v08_1x6.png}\\[3mm]
    \textbf{(c) 第 9 阶}\\[1mm]
    \includegraphics[width=\linewidth]{combined/layout_1x6_v9/FSCS_v09_1x6.png}\\[3mm]
    \textbf{(d) 第 10 阶}\\[1mm]
    \includegraphics[width=\linewidth]{combined/layout_1x6_v9/FSCS_v10_1x6.png}\\[3mm]
    \textbf{(e) 第 11 阶}\\[1mm]
    \includegraphics[width=\linewidth]{combined/layout_1x6_v9/FSCS_v11_1x6.png}\\[3mm]
    \textbf{(f) 第 12 阶}\\[1mm]
    \includegraphics[width=\linewidth]{combined/layout_1x6_v9/FSCS_v12_1x6.png}\\[3mm]
    \caption{自由-简支-固支-简支（FSCS）方板第 7 至第 12 阶屈曲模态临界荷载系数对比}
    \label{CH4_buckling_FSCS_b}
\end{figure}

% TODO[图像核对] 补充本段观察句：自由边工况下标准解模式（标准解面板的振型）与其他方法面板的一致性，
% 以及 MF 与 FEM 在自由边附近的挠度分布差异。

自由-简支-简支-简支（FSSS）工况含一个自由边，其屈曲临界荷载系数在六种工况中最低，结果如图\ref{CH4_buckling_FSSS_a} 与图\ref{CH4_buckling_FSSS_b} 所示。

\begin{figure}[H]
    \centering
    \textbf{(a) 第 1 阶}\\[1mm]
    \includegraphics[width=\linewidth]{combined/layout_1x6_v9/FSSS_v01_1x6.png}\\[3mm]
    \textbf{(b) 第 2 阶}\\[1mm]
    \includegraphics[width=\linewidth]{combined/layout_1x6_v9/FSSS_v02_1x6.png}\\[3mm]
    \textbf{(c) 第 3 阶}\\[1mm]
    \includegraphics[width=\linewidth]{combined/layout_1x6_v9/FSSS_v03_1x6.png}\\[3mm]
    \textbf{(d) 第 4 阶}\\[1mm]
    \includegraphics[width=\linewidth]{combined/layout_1x6_v9/FSSS_v04_1x6.png}\\[3mm]
    \textbf{(e) 第 5 阶}\\[1mm]
    \includegraphics[width=\linewidth]{combined/layout_1x6_v9/FSSS_v05_1x6.png}\\[3mm]
    \textbf{(f) 第 6 阶}\\[1mm]
    \includegraphics[width=\linewidth]{combined/layout_1x6_v9/FSSS_v06_1x6.png}\\[3mm]
    \caption{自由-简支-简支-简支（FSSS）方板第 1 至第 6 阶屈曲模态临界荷载系数对比}
    \label{CH4_buckling_FSSS_a}
\end{figure}
\begin{figure}[H]
    \centering
    \textbf{(a) 第 7 阶}\\[1mm]
    \includegraphics[width=\linewidth]{combined/layout_1x6_v9/FSSS_v07_1x6.png}\\[3mm]
    \textbf{(b) 第 8 阶}\\[1mm]
    \includegraphics[width=\linewidth]{combined/layout_1x6_v9/FSSS_v08_1x6.png}\\[3mm]
    \textbf{(c) 第 9 阶}\\[1mm]
    \includegraphics[width=\linewidth]{combined/layout_1x6_v9/FSSS_v09_1x6.png}\\[3mm]
    \textbf{(d) 第 10 阶}\\[1mm]
    \includegraphics[width=\linewidth]{combined/layout_1x6_v9/FSSS_v10_1x6.png}\\[3mm]
    \textbf{(e) 第 11 阶}\\[1mm]
    \includegraphics[width=\linewidth]{combined/layout_1x6_v9/FSSS_v11_1x6.png}\\[3mm]
    \textbf{(f) 第 12 阶}\\[1mm]
    \includegraphics[width=\linewidth]{combined/layout_1x6_v9/FSSS_v12_1x6.png}\\[3mm]
    \caption{自由-简支-简支-简支（FSSS）方板第 7 至第 12 阶屈曲模态临界荷载系数对比}
    \label{CH4_buckling_FSSS_b}
\end{figure}

% TODO[图像核对] 补充本段观察句：自由度约束最弱的 FSSS 工况下 MF 方法相对 FEM 的精度优势是否保持，
% 以及三种边界条件（FSCS、FSSS）之间的规律是否一致。

最后计算简支-固支交替（SCSC）工况。单向压缩沿 $x$ 方向，SCSC 与 CSCS 的边编号相差一个 $90^\circ$ 旋转而荷载方向不变，两种工况的屈曲临界荷载系数并不相等，结果如图\ref{CH4_buckling_SCSC_a} 与图\ref{CH4_buckling_SCSC_b} 所示。

\begin{figure}[H]
    \centering
    \textbf{(a) 第 1 阶}\\[1mm]
    \includegraphics[width=\linewidth]{combined/layout_1x6_v9/SCSC_v01_1x6.png}\\[3mm]
    \textbf{(b) 第 2 阶}\\[1mm]
    \includegraphics[width=\linewidth]{combined/layout_1x6_v9/SCSC_v02_1x6.png}\\[3mm]
    \textbf{(c) 第 3 阶}\\[1mm]
    \includegraphics[width=\linewidth]{combined/layout_1x6_v9/SCSC_v03_1x6.png}\\[3mm]
    \textbf{(d) 第 4 阶}\\[1mm]
    \includegraphics[width=\linewidth]{combined/layout_1x6_v9/SCSC_v04_1x6.png}\\[3mm]
    \textbf{(e) 第 5 阶}\\[1mm]
    \includegraphics[width=\linewidth]{combined/layout_1x6_v9/SCSC_v05_1x6.png}\\[3mm]
    \textbf{(f) 第 6 阶}\\[1mm]
    \includegraphics[width=\linewidth]{combined/layout_1x6_v9/SCSC_v06_1x6.png}\\[3mm]
    \caption{简支-固支交替（SCSC）方板第 1 至第 6 阶屈曲模态临界荷载系数对比}
    \label{CH4_buckling_SCSC_a}
\end{figure}
\begin{figure}[H]
    \centering
    \textbf{(a) 第 7 阶}\\[1mm]
    \includegraphics[width=\linewidth]{combined/layout_1x6_v9/SCSC_v07_1x6.png}\\[3mm]
    \textbf{(b) 第 8 阶}\\[1mm]
    \includegraphics[width=\linewidth]{combined/layout_1x6_v9/SCSC_v08_1x6.png}\\[3mm]
    \textbf{(c) 第 9 阶}\\[1mm]
    \includegraphics[width=\linewidth]{combined/layout_1x6_v9/SCSC_v09_1x6.png}\\[3mm]
    \textbf{(d) 第 10 阶}\\[1mm]
    \includegraphics[width=\linewidth]{combined/layout_1x6_v9/SCSC_v10_1x6.png}\\[3mm]
    \textbf{(e) 第 11 阶}\\[1mm]
    \includegraphics[width=\linewidth]{combined/layout_1x6_v9/SCSC_v11_1x6.png}\\[3mm]
    \textbf{(f) 第 12 阶}\\[1mm]
    \includegraphics[width=\linewidth]{combined/layout_1x6_v9/SCSC_v12_1x6.png}\\[3mm]
    \caption{简支-固支交替（SCSC）方板第 7 至第 12 阶屈曲模态临界荷载系数对比}
    \label{CH4_buckling_SCSC_b}
\end{figure}

% TODO[图像核对] 补充本段观察句：SCSC 与 CSCS 两种工况的对比结果是否互为对应，
% 以及 SCSC 工况下各方法对第 1 阶模态的临界荷载系数差异。

\section{小结}

% TODO: 第四章小结

\chapter{结论与展望}

% TODO: 待补充——全文结论（对应第3、4章数值结果）与研究展望
% 展望可预留一句：斜板及各向异性板可作进一步研究（2026-09-12 用户决定斜板暂不进正文）

% \include{acknowledgments}
\appendix
\chapter{应变分量逐项展开过程}
\label{app:strain-expansion}

% TODO: 待补充附录内容（式(4.4)面内应变按 x_3 幂次分项、式(4.5)横向剪切与法向应变的逐项代入过程）

\backmatter

\bibliography{references}

\include{cv}
\include{ai}

\end{document}
